%midterm2.tex
%%%%%%%%%%%%%%%%%%%%%%%%%%%%%%%%%%%%%%%%%%%%%%%%%%%%%%%%%%%%%%%%%%%%%%
%%%%%%            Math 403 Midterm \#2, Spring 2025             %%%%%%
%%%%%%                                                          %%%%%%
%%%%%%                 Instructor: Ezra Miller                  %%%%%%
%%%%%%%%%%%%%%%%%%%%%%%%%%%%%%%%%%%%%%%%%%%%%%%%%%%%%%%%%%%%%%%%%%%%%%
\documentclass[11pt]{article}

\oddsidemargin=17pt \evensidemargin=17pt
\headheight=9pt     \topmargin=26pt
\textheight=564pt   \textwidth=436.5pt
\parskip=1ex        %voffset=-6ex

\usepackage{amsmath,amssymb,graphicx,color,enumitem,amsthm}%,setspace%,mathrsfs%,hyperref
\usepackage{tikz-cd}

% Define theorem environment
\newtheorem{theorem}{Theorem}
\setlist[enumerate]{parsep=1.5ex plus 4pt,topsep=0ex plus 4pt,itemsep=0ex}
  \setlist[itemize]{parsep=1.5ex plus 4pt,topsep=-1ex plus 4pt,itemsep=-1.33ex}

\leftmargini=5.5ex
\leftmarginii=3.5ex

%for \marginpar to fit optimally
%hoffset=-1.02in
\setlength\marginparwidth{2.2in}
\setlength\marginparsep{1mm}
\newcommand\red[1]{\marginpar{\linespread{.85}\sf%
		   \vspace{-1.4ex}\footnotesize{\color{red}#1}}}
\newcommand\score[1]{\marginpar{\vspace{-2ex}\color{blue}{#1/14}}\hspace{-1ex}}
\newcommand\extra[1]{\marginpar{\color{blue}{#1/1}}\hspace{-1ex}}
\newcommand\total[2]{\marginpar{\colorbox{yellow}{\huge #1/#2}}}
\newcommand\collab[1]{\marginpar{\vspace{-11ex}\colorbox{yellow}{#1/3}}\hspace{-1ex}}
\newcommand\magenta[1]{\colorbox{magenta}{\!\!#1\!\!}}
\newcommand\yellow[1]{\colorbox{yellow}{\!\!#1\!\!}}
\newcommand\green[1]{\colorbox{green}{\!\!#1\!\!}}
\newcommand\cyan[1]{\colorbox{cyan}{\!\!#1\!\!}}
\newcommand\rmagenta[2][0ex]{\red{\vspace{#1}\magenta{\phantom{:}}\,: #2}}
\newcommand\ryellow[2][0ex]{\red{\vspace{#1}\yellow{\phantom{:}}\,: #2}}
\newcommand\rgreen[2][0ex]{\red{\vspace{#1}\green{\phantom{:}}\,: #2}}
\newcommand\rcyan[2][0ex]{\red{\vspace{#1}\cyan{\phantom{:}}\,: #2}}

\newcommand\points{\makebox[0pt][r]{[14pts]\ }}

%For separated lists with consecutive numbering
\newcounter{separated}

\newcommand{\excise}[1]{}
\newcommand{\comment}[1]{{$\star$\sf\textbf{#1}$\star$}}

%%%%%%%%%%%%%%%%%%%%%%%%%%%%%%%%%%%%%%%%%%%%%%%%%%%
%                                                 %
% TO ENABLE GRADING, DO NOT ALTER ABOVE THIS LINE %
%                                                 %
%%%%%%%%%%%%%%%%%%%%%%%%%%%%%%%%%%%%%%%%%%%%%%%%%%%

%new math symbols taking no arguments
\newcommand\0{\mathbf{0}}
\newcommand\bb{\mathbf{b}}
\newcommand\ee{\mathbf{e}}
\newcommand\pp{\mathbf{p}}
\newcommand\vv{\mathbf{v}}
\newcommand\ww{\mathbf{w}}
\newcommand\xx{\mathbf{x}}
\newcommand\yy{\mathbf{y}}
\newcommand\zz{\mathbf{z}}
\newcommand\CC{\mathbb{C}}
\newcommand\FF{\mathbb{F}}
\newcommand\RR{\mathbb{R}}
\newcommand\ZZ{\mathbb{Z}}
\newcommand\Cb{\mathbf{C}}
\newcommand\Nb{\mathbf{N}}
\newcommand\Ec{\mathcal{E}}
\newcommand\Mc{\mathcal{M}}
\newcommand\Pc{\mathcal{P}}
\newcommand\Vc{\mathcal{V}}
\newcommand\GL{\mathit{GL}}
\newcommand\from{\leftarrow}
\newcommand\ffrom{\longleftarrow}
\newcommand\goesto{\rightsquigarrow}

%redefined math symbols taking no arguments
\newcommand\<{\langle}
\renewcommand\>{\rangle}
\renewcommand\aa{\mathbf{a}}
\renewcommand\iff{\Leftrightarrow}
\renewcommand\implies{\Rightarrow}

%to explain about LaTeX commands:
\newcommand\command[1]{\texttt{$\backslash$#1}}

%for easy 2 x 2 matrices
\newcommand\twobytwo[1]{\left[\begin{array}{@{}cc@{}}#1\end{array}\right]}

%for easy column vectors
\newcommand\colvec[1]{\left[\begin{array}{@{}c@{}}#1\end{array}\right]}

%replaces \atop
\newcommand{\aoverb}[2]{{\genfrac{}{}{0pt}{1}{#1}{#2}}}

%for overlines
\def\ol#1{{\overline {#1}}}


%%%%%%%%%%%%%%%%%%%%%%%%%%%%%%%%%%%%%%%%%%%%%%%%%%%%%%%%%%%%%%%%%%%%%%
\begin{document}%%%%%%%%%%%%%%%%%%%%%%%%%%%%%%%%%%%%%%%%%%%%%%%%%%%%%%
%%%%%%%%%%%%%%%%%%%%%%%%%%%%%%%%%%%%%%%%%%%%%%%%%%%%%%%%%%%%%%%%%%%%%%


\title{\mbox{}\\[-8ex]Math 403 Midterm \#2, Spring 2025\\\normalsize
	Instructor: Ezra Miller\\[-1.5ex]}
\author{[1pt] Solutions by: Thomas Kidane\\[1ex]
I verify that I have consulted no electronic or human sources except
those allowed\\
and that I have in all other ways complied with Duke's code of
academic honesty.\hfill\mbox{}\\[1ex]
\ \ \ [1pt] Signature : I agree}
\date{Due Date: 11:59pm Thursday 17 April 2025}

\maketitle\total{}{100}
\thispagestyle{empty}%

\vspace{-3ex}%
\noindent
\textsc{Problems}

\begin{enumerate}
\item\ \score{}%1.
Fix an $n \times n$ complex matrix~$A$.  Among all of the eigenvalues
of~$A$, let $\lambda$ be one of maximal absolute value.  If $\sigma$
is the largest singular value of~$A$, show that $|\lambda| \leq
\sigma$.

\textbf{Solution}

It is enough to look at points on the unit circle/sphere/hyper-sphere, since all 
other vectors are just linear scalings of these vectors. Take the shmidt decomposition of 
of $A$.

\begin{align*}
	Ax=\sum_{k=1}^{n}{\sigma_{k}(x,v_{k})w_{k}}
\end{align*}

Where $n$ is the number of singular values(for the sake of simplicty I also consider singular values that are 0), the $v$'s and $w$'s orthogonal.
If we show that the norm of $Ax$ for any unit vector $x$ is bounded by $\sigma_{max}$, then 
any vector vector in any eigenspace is bounded by $\sigma$ and therfore $\lambda$ is also bounded by $\sigma_{max}$.

Take any $x$ in the unit sphere, then it can be decomposed into $v$'s, because the $v$'s are an othonormal basis.

\begin{align*}
	x=\alpha_{1}v_{1}+\cdots +\alpha_{n}v_{n}\\
	|x|^{2}=1=\alpha_{1}^{2}|v_{1}|^{2} +\cdots \alpha_{n}^{2}|v_{n}|^{2}\\
	1=\alpha_{1}^{2} +\cdots \alpha_{n}^{2}
\end{align*}

There are no mixed terms since the $v$'s are orthogonal.

\begin{align*}
	|A(\alpha_{1}v_{1}+\cdots+\alpha_{n}v_{n})|^{2}&=|\alpha_{1}w_{1}\sigma_{max}+\cdots +\sigma_{n}\alpha_{n}\sigma_{n}|\\
	&\alpha_{1}^{2}\sigma_{max}+\cdots +\alpha_{n}^{2}\sigma_{n}\leq \sigma_{max}
\end{align*}

There are no mixed terms since the $w$'s are orthogonal. Since the magnitude of the unit vectors are 
bounded by the maximum singular value, assume some $\lambda$ is the max eigenvalue. Then take the unit vector,
we know that this is bounded by $\sigma_{max}$ after applying $A$ to it,therefore $\lambda\leq \sigma_{max}$.

\item\ \score{}%2.
Fix $A \in \CC^{n \times n}$ and $\varepsilon > 0$.  Prove that there
is a consistent matrix norm $\Vert\cdot\Vert$ such that $\rho(A) \leq
\Vert A \Vert \leq \rho(A) + \varepsilon$.  Conclude that the spectral
radius of~$A$ satisfies $\rho(A) = \inf_{\Vert\cdot\Vert}\Vert A
\Vert$, where the infimum is taken over all consistent matrix norms
$\Vert\cdot\Vert$.

\textbf{Solution}

Lemma 1: The infinity norm is a consistent norm. Refer to Matrix Perturbation theory, 
chapter 2. 

Lemma 2: $v_{T}(X):=v(TXT^{-1})$ is consistent norm if $v$ is a consistent norm and $T$ is invertible.

Pf: We want $v_{T}(AB)\leq v_{T}(A)v_{T}(B)$, but if we plug in our definition, we 

\begin{align*}
	v(TABT^{-1}) \leq v(TAT^{-1})v(TBT^{-1})
\end{align*}

But this is true since $v$ is consistent(Take the matrix being applied to as $A=TAT^{-1}$ and $B=TBT^{-1}$).

Now we can take the jordan normal form of $A$, then scale the basis so that the elements 
on the super diagonal are less than $\epsilon$. Then we can make a consistent norm that applies the infinity norm 
on this(We can use lemma 2 to get rid of the unitary matrices in the jordan normal form decomposition). And we are done, since 
the maximum sum of the absolute values of the elements in the rows are less than $\rho(A)+\epsilon$. We can make a 
sequence of decreasing norms, but the lower bound is $\rho(A)$ since for any consistent norm $\rho(A)\leq v(A)$, and we can make 
a norm such that $v(A)<\rho(A) +\epsilon $ for any $\epsilon >0$. 

\item\ \score{}%3.
Show that equality in Elsner's theorem happens if and only if $A =
\omega\Vert A\Vert I$ for some complex number~$\omega$ with $|\omega|
= 1$ and $\widetilde A$ has $-\omega\Vert \widetilde A\Vert$ as an
eigenvalue, where $\Vert A\Vert$ is the operator norm of~$A$.

\textbf{Solution}

For the take of simplicity I have listed the original proof for Elsner's theorem below so that refer to it line by line, credit to G.W. Stewart and Ji-guang Sun for their book.

\begin{theorem}[1.3 (Elsner)]
For any \( A \) and \( \tilde{A} \),

\[
\text{hd}(A, \tilde{A}) \leq (\|A\|_2 + \|\tilde{A}\|_2)^{1 - \frac{1}{n}} \|E\|_{2}^{\frac{1}{n}}.
\]
\end{theorem}

\begin{proof}
Since the right-hand side of (1.4) is symmetric in \( A \) and \( \tilde{A} \), it is sufficient to prove that it bounds \( \text{sv}_A(\tilde{A}) \). Assume the maximum in (1.1) is attained for the eigenvalue \( \lambda \) of \( \tilde{A} \), and let \( x_1, \ldots, x_n \) be orthonormal vectors with \( \tilde{A}x_1 = \lambda x_1 \). Then
\[
\begin{aligned}
\text{sv}_A(\tilde{A})^n &\leq \Pi_i |\lambda_i - \tilde{\lambda}| \\
&= \det(A - \tilde{\lambda}I) \\
&\leq \Pi_i \|(A - \tilde{\lambda}I)x_i\|_2 \quad [\text{Hadamard inequality (1.2.2)}] \\
&= \|(A - \tilde{A})x_1\|_2 \Pi_{i>1} \|(A - \tilde{\lambda}I)x_i\|_2 \\
&\leq \|E\|_2 (\|A\|_2 + \|\tilde{A}\|_2)^{n-1}.
\end{aligned}
\]
\end{proof}

Suffieciecy($\rightarrow$):

Since the eigenvalues of A are the same we all only need to focus on the eigenvalues of $\tilde A$, and since it has the eigenvalue,
$-\omega \tilde A$ which is in the opposite direciton of $\omega$ the biggest difference is between this two eigenvalues, because 
any other eigenvalue can only be closer and not farther since this eigenvalue is in the opposite direciton and has the largest magnitude possible. 

Therefore $sv_{\tilde A}^{n}=( ||\omega A|| + || \omega \tilde A||)^{n} = ( || A|| + || \tilde A||)^{n} $(you can just add the magnitudes to get the difference), we also know that $||E||$ is $ ( || A|| + || \tilde A||)$ from the last line in the 
proof of Elsner's Theorem. We also know that $||((A - \tilde \lambda I))x_{i}||=|| \omega A|| +  || \omega \tilde A||=( || A|| + || \tilde A||)$. So we have equality.

Necessity($\leftarrow$):

For the first line to become equality we need that all the eigenvalues of $A$ have to be the same, let this value be $\omega ||A||$(The eigenvalue's magnitude will be the operator norm if all the eigenvalues are the same) . For the hadamard inequality to become equality on the third line we need that 
$||(A-\tilde \lambda)x_{i}||_{2}$ are an orthogonal set(That is when the Hadamard inequality becomes equality). Continuing on the transition from the $4^{th}$ to last line we need to have that,
$||(A-\tilde \lambda)x_{i}||_{2}=(||A||_{2}+||\tilde A||_{2})$, but for the triangle inequality to become equality, we need that the vectors be parallel.
So the $x_{i}$'s are an eigenvector of $A$. I couldn't prove that $A$ is diagonal(it would be diagonal if all vectors where in its eigenspace but I can't guarantee such a strong condition, since it depends $\tilde A$), but I can still say that say that $\tilde \lambda$ would have to be 
$-\omega ||\tilde A||$, since it has to be in the opposite direciton of $\omega ||A||$. Ofcourse $|\omega|$ has to be 1, since the eigenvalues of $|A|$ are all the same.
And from this we also know that $||E||=||A||+||\tilde A||$. Completing the proof. On a sidenote, $A$ would have to be diagonal, if it was written the basis of the eigenvectors of $\tilde A$, then it would satisfy the property, but I can't guarantee that either.

\item\ \score{}%4.
Prove that the coefficients of the characteristic polynomial
$p_\varphi$ of a linear transformation $\varphi: V \to V$ of a vector
space $V$ of dimension~$n$ all have (basis-free!)\ interpretations in
terms of determinants and traces by showing that $p_\varphi(t) =
\sum_j (-1)^j \mathrm{tr}(\bigwedge^j \varphi)\,t^{n-j}$.

\textbf{Solution}

Assume we are working in basis of eigenvectors. 

\begin{align}
	\varphi(v_{i})&=\lambda_{i}v_{i}\\
	tr\left( \bigwedge^r \varphi\right)&= \sum_{\sigma \in \binom{[dim V]}{r}}\bigwedge_{i=1}^{r}\varphi(v_{\sigma_{i}})\\
	tr\left( \bigwedge^r \varphi\right)&= \sum_{\sigma \in \binom{[dim V]}{r}}\prod_{i=1}^{r}\lambda_{\sigma_{i}}\\
	p_{\varphi}(t)&=(t-\lambda_{1})*\cdots *(t-\lambda_{n})\\
	p_{\varphi}(t)&=\sum_{i=0}^{dim V}(-1)^{dimV-i}t^{i}\sum_{\sigma \in \binom{[dim V]}{r}}\prod_{i=1}^{r}\lambda_{\sigma_{i}}\\
	p_{\varphi}(t)&=\sum_{i=0}^{dim V}(-1)^{dimV-i}t^{i}\sum_{\sigma \in \binom{[dim V]}{r}}\bigwedge_{i=1}^{r}\varphi(v_{\sigma_{i}})\\
	p_{\varphi}(t)&=\sum_{i=0}^{dim V}(-1)^{dimV-i}tr\left( \bigwedge^r \varphi\right)t^{i}\\
\end{align}

We extend the proof to any basis(basis-free) by noting that $\bigwedge \varphi$ is a linear function, therefore the 
trace doesn't change since it is a property of the linear transformatiom, not its representation. So the trace is the same in any basis,
therefore this fact works independent of the basis you working in.

\item\ \score{}%5.
Let $A \in \RR^{2 \times 2}$ be a diagonal matrix.  What conditions
on~$A$ guarantee that the exponential map $\exp_{iA}: \RR \to
\GL_2(\CC)$ sending $t \mapsto e^{itA}$ is injective?  Set $G =
\exp_{iA}(\RR)$, the image of $\exp_{iA}$.  When $\exp_{iA}$ is
injective, can it be that for every open interval $I \subseteq \RR$
there is an open subset $U \subseteq \GL_2(\CC)$ such that
$\exp_{iA}(I) = U \cap G$?

\textbf{Solution}

Since $A$ is diagonal that means
\begin{align*}
	A=\begin{bmatrix}
		a & 0\\
		0 & b
	\end{bmatrix}\\
	e^{itA}=\begin{bmatrix}
		e^{ita} & 0\\         	
		0 & e^{itb}
	\end{bmatrix}
\end{align*}

For the function to be injective it needs to be the case,
$exp_{iA}(t_{1}) \neq exp_{iA}(t_{2}) \Longleftrightarrow t_{1}\neq t_{2}$. 

$e^{ita}$ goes in a cycle of period $\frac{2\pi}{a}$, same thing for $e^{itb}$, with period $\frac{2\pi}{b}$.

If $\exists t_{1}$ and $t_{2}$ s.t $t_{1} \neq t_{2}$ and $exp_{iA}(t_{1})=exp_{iA}(t_{2})$, then that means the 
two values must have cycled, which means that you can find $x,y\in \ZZ$ s.t $\frac{2x\pi}{a}=\frac{2y\pi}{b} \rightarrow xb=ay\rightarrow a/b\in \mathbb{Q}$.
So $a/b$ must be irrational for the map to be injective, this condition is also sufficient, since for the two values to repeat, it must be 
that case that the ratio of their periods rational otherwise you couldnt' find a period in which both the values would repeat.

I think the second value is false. Imagine $\RR^8$(very hard I know), this space is close enough(there are some matrices that aren't complex linear but this proof lost its rigour already) to the space of 
invertible complex matrices. The question is equivalent to asking if I can make an open subset in this subspace,
s.t its intersection with the image of the function is some range. Assume there is some open set you can make such that this is true. I argue, 
that there is a point in that subspace that is not in your range, my justification is that I can make an arbitrarily close approximation of any value in your range,
but its inverse will be arbitrarily far.

Let's take a concrete example, say $a=1$ and $b=\sqrt{2}$. Assume you made some open subset in $\GL_2(\CC)$, such the 
intersection of this subspace with the image corresponded to the interval (-1,1). Obviously, $\GL_2(\CC)$ has a norm, you can 
view it as a space that is isomorphic to $\RR^{8}$ and do the normal euclidean norm. I argue I can find a matrix in the image, that is 
arbitrarily close to $I$ in euclidean distance terms, but not in the interval because I can find some arbitrarily large value of $x$
such that $|I-exp_{iA}(x)|<\epsilon$, for any $\epsilon>0$. Why? assume I fix $\epsilon$, then I keep increasing $x$, the values of $e^{ita}$ and $e^{itb}$ keep looping, but they are never in sync,
everytime there is some grap between their positions, this gap never goes to zero, but it doesn't diverge either, it keeps oscillating, and at some point with 
probabilty $1^{*}$, it gets close enough to 0, arbitrarily closely(therefore arbitrarily close to $I$).

The basic idea of my proof is that, the image is too dense, any level of precision is insuffient since there are
infinitely many ranges in a nonzero amount of volume(you have fit the whole number line in the small hypercube).

\item\ \score{}%6.
Fix $A \geq 0$ in $\RR^{n \times n}$.  Let $a^{(k)}_{ij}$ be the $ij$
entry of $A^k$.  Assume for each $i$ and~$j$ that~$a^{(k)}_{ij} \neq
0$ for some $k \in \ZZ_{> 0}$ that may depend on~$i$ and~$j$.  Prove
that $(I_n + A)^{n-1} > 0$.

\textbf{Solution}

The exact magnitudes of the entries in the $A$ don't matter, replace all positve entries with 1 from now on.
Imagine $n$ cities, let there be a one directional road from city $i$ to city $j$, if $a_{ij}=1$(We replaced the positive entries
of $A$ with 1, so this is equivalent to saying that the original entry had a positive value).

Taking the powers of $A$ is equivalent to taking walkes along this city(graph), more specifically
$a_{ij}^{(k)}>0$ tells you whether there is a path that consists of $k$ roads from city $i$ to city $j$. 
Therefore, saying that $a_{ij}^{(k)}>0$ for some $k$ is equivalent to saying that a path exists from city $i$ to city
$j$ which consists of $k$ roads for some $k$.

\begin{align*}
	(I+A)^{n-1}&=\sum_{i=0}^{n-1}\binom{n-1}{i} I^{i}A^{n-1-i}\\
	&= I + (n-1)A + \cdots + A^{n-1}
\end{align*}

For an entry in $(I+A)^{n-1}$ to be zero it must mean that there is no path between some two cities that containts
any number of roads between 0 to $n-1$ roads, but that is just nonsense because we know that there is a 
path between any two citeis and no path between any two cities need not be longer that $n-1$(if you start at one city and can get to another city, then 
the shortest path between the city you are in and the other city, can't be longer than $n-1$, otherwise you had to revist one city). $(I+A)^{n-1}$ effectively 
takes a snapshot of all paths between any two cities, and since we said that you can get from one city to any other city after some number of walks, all the 
entries in $(I+A)^{n-1}$ must be strictly positive.

\item\ \score{}%7.
Suppose that $U$ is a normal operator whose eigenvalues all lie on the
unit circle in~$\CC$.  Must $U$ be unitary?  Give a proof or
counterexample.

\textbf{Solution}

From Theorem 2.4 in Chapter 6 Lax, we know that any normal operator $U$ in a complex vector space has an orthonormal basis of 
eigenvectors. Take the decomposition of $U$ to be $PDP^{*}$, where $P$ is unitary and $D$ diagonal. If we  multiply this by its conjuguate transpose. $PDP^{*}PD^{*}P^{*}=PD^{*}DP^{*}=PDD^{*}P^{*}\rightarrow DD^{*}=D^{*}D\rightarrow I$,
since the values on the diagonal have absolute value 1, and multiplying a complex number by its conjugate gets you its magnitude. Therefore $N^{*}$ is the inverse of $N$, since $PIP^{*}=I$, since $P$ is 
unitary, which means $N$ is unitary.



\end{enumerate}
\end{document}


%%% Various Emacs customizations:
%%% Local Variables:
%%% fill-column: 70
%%% indent-tabs-mode: t
%%% TeX-electric-sub-and-superscript: nil
%%% TeX-brace-indent-level: 0
%%% LaTeX-indent-level: 0
%%% LaTeX-item-indent: 0
%%% End:
